\begin{definition}
For a two-bite sample \(\omega\) and a base vertex \(x\), the finite set
\(\noderef{TwoBiteLiftedPositiveNeighborhood}(\omega,x)\subseteq[n]\) is
defined by the following two cases.  If \(x=\operatorname{red}(r)\), then a
vertex \(v\in[n]\) belongs to the set exactly when
\[
r<\pi_R^\omega(v)
\quad\text{and}\quad
\noderef{TwoBiteRedGraph}(\omega)\text{ contains the red base edge }
\{r,\pi_R^\omega(v)\}.
\]
If \(x=\operatorname{blue}(b)\), then \(v\) belongs exactly when
\[
b<\pi_B^\omega(v)
\quad\text{and}\quad
\noderef{TwoBiteBlueGraph}(\omega)\text{ contains the blue base edge }
\{b,\pi_B^\omega(v)\}.
\]
Equivalently, this is the union of the fibers over same-color neighbors of
\(x\) with strictly larger base index.  These are the one-sided neighborhoods
used by the paper when a monochromatic triangle is charged to its
lexicographically latest same-color edge.
\end{definition}
